\chapter{The Big Questions and a few smaller Ones}

\section*{Can computer feel pain in principle?}

A computer, even a simple thermostat, given sufficient time, could in theory duplicate the informational state of a suffering human.
However, this does not make the device itself conscious. There is no residual mystery in the physical operation of the machine.

\section{Can computer simulate pain?}

We  observers are static finite set of information. We cannot access new information - just explore the information in our set. 
One of these explorations could take the form of DNA simulation. It therefore cannot create new suffering Alice.
It discovers (traces out one) of the very many configurations that describe an equivalence class - Alice.
She exists regardless of whether a human-built machine ever runs her code.


\section*{The source of the Power of Quantum Computers}

Some theorists, like David Deutsch \cite{deutsch1997fabric}, argue that quantum computers achieve their speed by performing calculations across many parallel universes simultaneously.
However, under the lens of this theory, quantum computers do not "create" new answers through labor.
They reveal information that already inherent exist. 

Much like the simulation of a human doesn't create new suffering but reveals a pre-existing configuration, a quantum computer uses interference to navigate a landscape of existing possibilities, revealing one of them into a single, observable result in our universe.


\section*{Is this a "One-Parameter" Theory rather than Zero?}

Strictly speaking, \(\lambda\) appears as a parameter in the probability measure: 

\[
\mathbb{P}(\gamma \mid O) \propto \exp[-\lambda\, \mathcal{C}_O(\gamma)]
\]

However, \(\lambda\) is not an arbitrary input like the mass of a Higgs boson.
It represents the Global Compression Pressure—the degree to which the informational substrate is "packed" into predictable patterns.

\paragraph{The Analogy of Surface Tension}
Think of \(\lambda\) like the surface tension of water.
One doesn't "need" to input the tension to have a drop of water; the tension is an emergent property of how the molecules interact to find their lowest-energy state. 

In this theory:
\begin{itemize}
    \item \(\lambda\) is the Selection Pressure: It determines how "strictly" the universe favors simplicity over chaos.
    \item If \(\lambda \to 0\), the universe is white noise (infinite complexity).
    \item If \(\lambda \to \infty\), the universe is a single, static bit (maximal simplicity).
\end{itemize}


\section*{Where did all the matter in the universe come from?}

This question implicitly assumes temporal ordering and an external origin.
Both time and causality are emergent, observer-dependent phenomena, as established in Chapter~\nameref{humans-as-axiomatic-systems}.


\section*{Why something rather than nothing?}

This is analogous to asking why ``heads'' rather than ``tails'': both are equally abstract and neither is privileged.

If there were a rule that favored nothingness over something, we would again have to ask about the origin of such a rule—who set it?
Non-empty structures are just as real or unreal as empty sets. Both exist because nothing forbids them.


\section*{Why do particles follow a complex-valued wavefunction?}

Because we are observing compressed structure.
The wavefunction is the \textbf{compression codec} of the underlying informational structure of reality.
Observed wave-like behavior emerges from this compression.

Formally:
\[
\boxed{\text{Maximal Predictability}} \;\rightarrow\; \boxed{\text{Maximal Compression}} \;\rightarrow\; \boxed{\text{Maximal Probability}} 
\]

Predictability is exactly what we call the laws of physics.

Configurations with minimal observer-indexed complexity \( \mathcal{C}_O[\gamma] \) dominate the probability measure:
\[
\mathbb{P}(\gamma \mid O) \propto \exp[-\lambda \, \mathcal{C}_O[\gamma]].
\]
Hence, the ``laws'' of physics are statistical consequences of overwhelmingly probable, compressible configurations.


\section*{Why does the universe expand? Is there a link to entropy?}

Spacetime expansion is the geometric manifestation of increasing configurational entropy. A bitstring of zero entropy corresponds to a singularity—a geometric mapping with zero volume. As the entropy of the configuration increases, the corresponding geometric representation must necessarily 'stretch' to accommodate the growing density of internal constraints.


\section*{Why did the universe start with zero entropy?}

Low-entropy configurations are maximally compressible and thus most probable.
High-entropy starting states would require additional information to encode the observer, increasing \( \mathcal{C}_O[\gamma] \) and suppressing probability:
\[
\mathbb{P}(\gamma_{\rm high-entropy} \mid O) \ll \mathbb{P}(\gamma_{\rm zero-entropy} \mid O).
\]

Consequently, the arrow of time and increasing entropy emerge naturally:
\[
\text{Zero entropy} \;\Rightarrow\; \text{Maximal compressibility} \;\Rightarrow\; \text{Observer-compatible universe}.
\]

Other observer-compatible configurations may exist with different geometric histories.


\section*{Is the Universe Infinite?}

If it were finite, one would have to ask: finite by how many bits? And who, or what, set that limit? Any fixed bound would itself require explanation.
A finite informational universe merely pushes the mystery back one level. 

So the only non-arbitrary conclusion is that the deep nature of everything cannot be finitely bounded.
At the most fundamental level, reality must be infinite—perhaps not even informational in the familiar sense, but beyond any finite description.

Based on all observational evidence, we observers however, are finite structures. According to this theory, observers
must be finite, because infinities don't compress well, which drops the induced measure.


\section*{What is gravity?}

We observe gravity for the same reason we observe waves - we are observing compressed structure.

Observers find themselves within the most compressible (probable) configurations, and the path through these configurations is what observers sense as time and space.
Each configuration is slightly different, which observers perceive as moving in geometric spacetime.

We ``fall'' because there are more copies of us ``down there.'' Gravity is the statistical bias toward the most probable (compressible) configurations describing the observer.


\section*{How does Gödel incompleteness affect the theory?}

Stephen Hawking (in “Gödel and the End of Physics”) argued that a physical Theory of Everything would be self-referencing (like a formal system describing itself), so by analogy with Gödel, it would likely be incomplete.
He, among many others, have suggested we may never have a finite, complete set of principles for the universe.

By placing $U$ outside formal systems and treating physics/math as emergent interpretive phenomenology, this framework largely dissolves the threat.
Incompleteness becomes expected (even welcome) at the symbolic layer, while the deeper ontological substrate, namely the interpretation space, remains untouched.

Within the present framework, formal systems are not fundamental features of reality, but emergent symbolic structures stabilized within observer equivalence classes.

The informational totality \(U\) itself is not a formal axiomatic system and therefore is not subject to Gödel incompleteness in any direct sense.


\section*{Are we living in a simulation?}

And who then would simulate the simulator?

We are living in a information. All observer-compatible configurations exist within a single informational substrate that is abstract by nature.


\section*{Does the theory support QBism?}

QBism interprets the wavefunction as a representation of an agent’s subjective degrees of belief, where observation is an active, participatory event in an unpredictable world. This framework completely rejects that view. 

Here, the universe is entirely static; there are no metaphysical signals, no temporal transitions, and no dynamic collapse.
The observer does not "discover" or "alter" reality through observation.
Instead, everything the observer perceives is information that already immutably coexists within her localized state. 

The wavefunction itself possesses no independent physical reality; it is entirely emergent.
Much like the way pixels in an MPEG-encoded movie appear to move and wave while the underlying data remains a frozen, compressed file on a disk, the quantum wavefunction is merely an artifact of compressed structure.
Probability and dynamics are not fundamental; they are the geometric illusions generated by decoding statistical typicality and compressibility within a timeless, abstract architecture.



\section*{Why do we age and die?}

The observer’s future is limited by the spectral complexity of their wavefunction, \(\psi_O\). Let \(\Sigma[\gamma]\) denote the number of independent frequency-phase components required to encode future history \(\gamma\).
The probability of persistence scales as:

\[
\mathbb{P}(\gamma \mid O) \propto \exp(-\alpha \, \Sigma[\gamma]).
\]

The accumulation of memories and internal structure increases \(\Sigma[\gamma]\), reducing the measure of compatible continuations.
Aging and death are the geometric and statistical manifestations of combinatorial exhaustion.


\section*{Is this just Solipsism?}

This theory takes the observer’s existence as axiomatic, but it does not imply solipsism. The number of observers is a probabilistic variable.
It would be very hard to argue that even from two different configurations; Alice with blond hair, and Alice with brunette hair,
only the other one would be allowed. One would have to show that a single observer is statistically favored.


\section*{What motivates the Spectral Complexity?}

Observational evidence. We observe that the micro-cosmos behaves as waves. 


\section*{Does the theory explain why the expansion rate is near-critical?}

Conditional on increasing entropy and typical emergence, observers are most likely to find themselves near the peak of the microstructure count (maximum compressibility).
This automatically corresponds to a near-critical expansion regime.


